% sec:kerr-horizons — Horizons, Ergosphere, and Frame Dragging
\section{Horizons, Ergosphere, and Frame Dragging}
\label{sec:kerr-horizons}

A non-rotating black hole has one important surface: the event
horizon at $r = 2M$.  Rotation splits this into two horizons and
creates a third surface --- the ergosphere boundary --- that has no
Schwarzschild analogue.  Inside the ergosphere, no observer can
remain stationary with respect to distant stars; the spacetime
itself drags everything in the direction of the hole's rotation.
Understanding these surfaces and the dragging they encode is
essential for the ray tracer: horizon radii set the termination
condition for captured photons, and the frame-dragging angular
velocity controls the brightness asymmetry between the approaching
and receding sides of the accretion disc.
\coderef{Spacetime.Kerr}{kerrHorizonRadius}

\paragraph{Event horizons.}

A horizon is a surface from which no outward-directed causal
signal can escape to infinity; in Boyer--Lindquist coordinates
this condition surfaces as a divergence of the radial metric
coefficient.  The horizons of the Kerr spacetime are the surfaces
where $g_{rr} = \Sigma/\Delta$ diverges, i.e.\ where the horizon
function $\Delta = r^2 - 2Mr + a^2$ from \cref{eq:kerr-Delta}
vanishes.  This is a quadratic in $r$ with roots
%
\begin{equation}
  r_\pm = M \pm \sqrt{M^2 - a^2} \,.
  \label{eq:kerr-horizons}
\end{equation}
%
Both roots are real and positive precisely when $\abs{a} \leq M$.
The outer root $r_+$ is the event horizon: no future-directed
causal curve can escape from $r < r_+$ to spatial infinity.  The
inner root $r_-$ is the Cauchy horizon, a surface beyond which the
initial-value problem for the field equations ceases to have a
unique solution.
\physnote{At the Cauchy horizon, signals from the singularity can
reach an infalling observer, so the past no longer uniquely
determines the future.  The physical relevance of the Kerr interior
below $r_-$ remains an open question in classical and quantum
gravity \cite{Wald1984}.}
Setting $a = 0$ recovers
$r_+ = 2M$ and $r_- = 0$, reproducing the Schwarzschild
horizon and singularity (\cref{sec:kerr-bl}).
\implnote{The codebase function \texttt{kerrHorizonRadius} computes
$r_+$ directly: \texttt{m + sqrt (m*m - a*a)}.  The geodesic
integrator terminates a ray when its Boyer--Lindquist radius drops
below $r_+$.}

For a $10\,M_\odot$ black hole, the event horizon shrinks from
$r_+ = 2M \approx 30\;\text{km}$ at zero spin to
$r_+ = M \approx 15\;\text{km}$ at the extremal limit $a = M$.
At extremality the two horizons merge:
$\Delta = (r - M)^2$, and the surface gravity vanishes,
giving a degenerate horizon whose thermodynamic
consequences are discussed in the literature
\cite{Bardeen1972}.

\paragraph{The ergosphere.}

A second characteristic surface arises from the time--time component
of the Boyer--Lindquist metric.  An observer at fixed
$(r, \theta, \varphi)$ has four-velocity proportional to
$\partial_t$, so the squared norm of the time Killing vector is
%
\begin{equation}
  g_{tt} = -\!\left(1 - \frac{2Mr}{\Sigma}\right)
         = -\frac{\Delta - a^2\sin^2\!\theta}{\Sigma} \,.
  \label{eq:kerr-gtt}
\end{equation}
%
The second form follows from
$\Sigma - 2Mr = (r^2 - 2Mr + a^2) - a^2\sin^2\!\theta
= \Delta - a^2\sin^2\!\theta$.
The quantity $g_{tt}$ vanishes when
%
\begin{equation}
  r = r_{\text{ergo}}(\theta) = M + \sqrt{M^2 - a^2\cos^2\!\theta} \,,
  \label{eq:kerr-ergosphere}
\end{equation}
%
which defines the outer boundary of the ergosphere (we discard the
smaller root, which lies inside the inner horizon).  The
discriminant $M^2 - a^2\cos^2\!\theta \geq 0$ follows from
$\abs{a} \leq M$, so $r_{\text{ergo}}$ is real for all
$\theta$.  The surface $r = r_{\text{ergo}}$ touches the event
horizon at the poles ($\theta = 0, \pi$), where
$r_{\text{ergo}} = r_+$, and bulges outward at the equator, where
$r_{\text{ergo}} = 2M$ regardless of the spin --- matching the
Schwarzschild event horizon radius.
\physnote{The region $r_+ < r < r_{\text{ergo}}$ is the
ergoregion.  Inside it, $g_{tt} > 0$: the Killing vector
$\partial_t$ is spacelike, and no observer can remain
stationary.  This is the arena for the Penrose process,
in which a particle entering the ergoregion splits so that
one fragment falls into the hole with negative energy (as
measured at infinity) while the other escapes with more
energy than the original --- extracting rotational energy
from the black hole \cite{Wald1984}.}

\begin{figure}[htbp]
  \centering
  % kerr-horizons-ergosphere.tex — Kerr horizons and ergosphere cross-section
% Inputable TikZ fragment for fig:kerr-horizons-ergosphere
% Requires: tikz with calc, arrows.meta (loaded by ehbook.cls)
%
% Meridional cross-section of the Kerr black hole structure for a/M = 0.9.
% Axes: horizontal = equatorial direction (r sin theta),
%        vertical  = spin axis (r cos theta).
% In BL coordinates, surfaces of constant r are circles (spheres in 3D).
% The ergosphere boundary r_ergo(theta) is oblate, touching r+ at poles.
%
% Parameters: M = 1, a/M = 0.9.
%   r+ = 1.4359M,  r- = 0.5641M,  r_ergo(equator) = 2M.
% Scale: 1.8 cm/M.
%   r+ radius = 2.585 cm,  r- radius = 1.015 cm,  r=2M radius = 3.600 cm.
%
\begin{tikzpicture}[
    >={Stealth[length=5pt]},
    font=\footnotesize,
  ]

  % ── Ergoregion shading (between ergosphere and r+) ──────────────────
  % Even-odd rule: fill between the ergosphere path and the r+ circle.
  \fill[EHAccretionGold, opacity=0.10, even odd rule]
    plot[smooth cycle, tension=0.7] coordinates {
      ( 0.000,  2.585)
      ( 0.457,  2.593)
      ( 0.944,  2.594)
      ( 1.464,  2.535)
      ( 1.995,  2.378)
      ( 2.504,  2.101)
      ( 2.951,  1.704)
      ( 3.301,  1.201)
      ( 3.524,  0.621)
      ( 3.600,  0.000)
      ( 3.524, -0.621)
      ( 3.301, -1.201)
      ( 2.951, -1.704)
      ( 2.504, -2.101)
      ( 1.995, -2.378)
      ( 1.464, -2.535)
      ( 0.944, -2.594)
      ( 0.457, -2.593)
      ( 0.000, -2.585)
      (-0.457, -2.593)
      (-0.944, -2.594)
      (-1.464, -2.535)
      (-1.995, -2.378)
      (-2.504, -2.101)
      (-2.951, -1.704)
      (-3.301, -1.201)
      (-3.524, -0.621)
      (-3.600,  0.000)
      (-3.524,  0.621)
      (-3.301,  1.201)
      (-2.951,  1.704)
      (-2.504,  2.101)
      (-1.995,  2.378)
      (-1.464,  2.535)
      (-0.944,  2.594)
      (-0.457,  2.593)
    }
    (0, 0) circle (2.585);

  % ── Schwarzschild reference: r = 2M ─────────────────────────────────
  \draw[EHMarginGray, dotted, thin]
    (0, 0) circle (3.600);
  \node[font=\scriptsize, text=EHMarginGray, anchor=south west]
    at (2.55, 2.55) {$r = 2M$};

  % ── Ergosphere boundary ─────────────────────────────────────────────
  \draw[EHAccretionGold, thick]
    plot[smooth cycle, tension=0.7] coordinates {
      ( 0.000,  2.585)
      ( 0.457,  2.593)
      ( 0.944,  2.594)
      ( 1.464,  2.535)
      ( 1.995,  2.378)
      ( 2.504,  2.101)
      ( 2.951,  1.704)
      ( 3.301,  1.201)
      ( 3.524,  0.621)
      ( 3.600,  0.000)
      ( 3.524, -0.621)
      ( 3.301, -1.201)
      ( 2.951, -1.704)
      ( 2.504, -2.101)
      ( 1.995, -2.378)
      ( 1.464, -2.535)
      ( 0.944, -2.594)
      ( 0.457, -2.593)
      ( 0.000, -2.585)
      (-0.457, -2.593)
      (-0.944, -2.594)
      (-1.464, -2.535)
      (-1.995, -2.378)
      (-2.504, -2.101)
      (-2.951, -1.704)
      (-3.301, -1.201)
      (-3.524, -0.621)
      (-3.600,  0.000)
      (-3.524,  0.621)
      (-3.301,  1.201)
      (-2.951,  1.704)
      (-2.504,  2.101)
      (-1.995,  2.378)
      (-1.464,  2.535)
      (-0.944,  2.594)
      (-0.457,  2.593)
    };

  % ── Outer horizon r+ ────────────────────────────────────────────────
  \draw[EHHorizonCrimson, thick]
    (0, 0) circle (2.585);

  % ── Inner horizon r- ────────────────────────────────────────────────
  \draw[EHHorizonCrimson, thick, dashed]
    (0, 0) circle (1.015);

  % ── Ring singularity ────────────────────────────────────────────────
  % At r = 0, theta = pi/2: equatorial axis at origin in BL cross-section.
  \fill[EHHorizonCrimson] (0, 0) circle (2pt);

  % ── Axes ────────────────────────────────────────────────────────────
  % Spin axis (vertical)
  \draw[->, EHMarginGray, thin]
    (0, -3.2) -- (0, 3.5)
    node[above, font=\scriptsize, text=EHMarginGray] {spin axis};
  % Equatorial direction (horizontal)
  \draw[->, EHMarginGray, thin]
    (-4.1, 0) -- (4.3, 0)
    node[right, font=\scriptsize, text=EHMarginGray] {equatorial plane};

  % ── Labels ──────────────────────────────────────────────────────────

  % r+ label with leader line
  \node[font=\scriptsize, text=EHHorizonCrimson, anchor=west]
    at (2.0, -2.3) {$r_+$};
  \draw[EHHorizonCrimson, thin, shorten >=1pt]
    (2.0, -2.3) -- (1.83, -1.83);

  % r- label with leader line
  \node[font=\scriptsize, text=EHHorizonCrimson, anchor=west]
    at (1.3, -0.85) {$r_-$};
  \draw[EHHorizonCrimson, thin, shorten >=1pt]
    (1.3, -0.85) -- (0.72, -0.72);

  % r_ergo label at equatorial bulge
  \node[font=\scriptsize, text=EHAccretionGold!85!black, anchor=west]
    at (3.75, 0.35) {$r_{\text{ergo}}$};

  % Ergoregion label
  \node[font=\scriptsize, text=EHAccretionGold!85!black, rotate=55]
    at (-2.3, 1.55) {ergoregion};

  % Ring singularity label
  \node[font=\scriptsize, text=EHHorizonCrimson, anchor=north west]
    at (0.12, -0.12) {$r = 0$};

\end{tikzpicture}

  \caption[Kerr horizons and ergosphere]{%
    Meridional cross-section of the Kerr black hole structure for spin
    $a/M = 0.9$, plotted in Boyer--Lindquist coordinates
    ($r\sin\theta$ horizontal, $r\cos\theta$ vertical).  The outer
    horizon $r_+$ (solid crimson) and inner horizon $r_-$ (dashed
    crimson) are circles because surfaces of constant BL radius are
    spheres.  The ergosphere boundary $r_{\text{ergo}}(\theta)$ (solid
    gold) touches $r_+$ at the poles and bulges out to $r = 2M$ at the
    equator, creating the ergoregion (gold shading).  The dotted grey
    circle marks $r = 2M$, the Schwarzschild event horizon radius.
    The ring singularity at $r = 0$, $\theta = \pi/2$ is marked at
    the origin.}
  \label{fig:kerr-horizons-ergosphere}
\end{figure}

\paragraph{Frame-dragging angular velocity.}

The off-diagonal $g_{t\varphi}$ term in the Boyer--Lindquist
metric~\cref{eq:kerr-bl-metric} couples time and azimuth, meaning
that a freely falling observer with zero conserved angular momentum
($p_\varphi = 0$) acquires azimuthal velocity.  The angular
velocity of such a zero-angular-momentum observer (ZAMO) is
%
\begin{equation}
  \omega = -\frac{g_{t\varphi}}{g_{\varphi\varphi}}
         = \frac{2Mar}{A} \,,
  \label{eq:kerr-omega}
\end{equation}
%
where $A = (r^2 + a^2)^2 - a^2\Delta\sin^2\!\theta$ is the
azimuthal auxiliary from \cref{eq:kerr-A}.  A nonzero $\omega$
means that a locally non-rotating observer is nonetheless carried
around the hole azimuthally, and photons inherit this rotation: the
spacetime drags everything in the direction of the black hole's
spin.  At the event horizon, $\Delta = 0$ gives
$A = (r_+^2 + a^2)^2$ and therefore
%
\begin{equation}
  \omega_H = \frac{2Mar_+}{(r_+^2 + a^2)^2}
           = \frac{a}{r_+^2 + a^2} \,,
  \label{eq:kerr-omega-horizon}
\end{equation}
%
using $2Mr_+ = r_+^2 + a^2$ (which follows from
$\Delta(r_+) = 0$).  This is the angular velocity of the
horizon itself: all objects, including photons, co-rotate
with the horizon at $r = r_+$.  For $a = 0$, $\omega_H = 0$
and the non-rotating limit is recovered.
\histnote{The angular velocity $\omega_H$ also appears in black
hole thermodynamics: the first law of black hole mechanics relates
changes in mass, area, and angular momentum through $\omega_H$
\cite{Bardeen1972}.}

At large radii, $A \approx r^4$ and $\omega$ falls off as
%
\begin{equation}
  \omega \approx \frac{2Ma}{r^3} \,,
  \label{eq:kerr-omega-far}
\end{equation}
%
the same $1/r^3$ fall-off as the angular velocity field of a
rotating body in the weak-field post-Newtonian approximation
(\cref{sec:gr-weak-field}).  For a
$10\,M_\odot$ black hole with $a/M = 0.9$ at $r = 10M$, this
gives $\omega \approx 1.8 \times 10^{-3}/M \approx 37\;\text{rad\,s}^{-1}$
--- fast enough to produce observable Doppler shifts in photons
from the accretion disc.

\paragraph{Frame dragging and the ray tracer.}

The ZAMO angular velocity enters the ray tracer indirectly through
the metric.  In Kerr--Schild coordinates, the frame-dragging
information is encoded in the twist of the null covector $l_\mu$
(\cref{eq:kerr-ks-null}): the terms $ay$ and $-ax$ in the
spatial components of $l_\mu$ mix the $x$ and $y$ directions by an
amount proportional to $a$.  The geodesic integrator does not
compute $\omega$ explicitly --- it integrates Hamilton's equations
using the full metric --- but the physical effect is the same:
photons co-rotating with the black hole follow shorter paths and
arrive blueshifted, while counter-rotating photons travel longer
paths and arrive redshifted.  This asymmetry is the dominant
visual feature in rendered images of Kerr black holes with
optically thin accretion flows.
\warnnote{Frame dragging is often described loosely as ``the
black hole dragging spacetime around with it''.  More precisely:
the Killing vector $\partial_t$ acquires a nonzero inner product
with $\partial_\varphi$, so the locally non-rotating frame
rotates relative to observers at infinity.  No physical medium is
being dragged.}
